\begin{abstract}
Case handlers must request evidence before a file settles the deciding facts. In \casepath, a model reads each such fact three times and must quote the file. Code admits a value only when the readings agree, executes the rules of a workflow, and holds every request that depends on an open fact. We tested a frozen commit on sealed synthetic claims under a pre-registered plan. On tenancy it made \sthreeHTwoA{} strict unsafe errors per claim, against \sthreeHTwoB{} for a single prompt with the same rules on the same cost-efficient model and \sthreeHThreeB{} on a frontier model, and held \sthreeCasePathTenCorrectHold\% of the documents the reference holds, against \sthreeRicTenCorrectHold\% and \sthreeRicFrontierTenCorrectHold\%. On theft it made \sthreeHOneA{} errors per pair against \sthreeHOneB{} for the same-model prompt. These three hypotheses were confirmed; its theft lead over self-consistency (\sthreeHFourSCB) missed significance, so a fourth failed and a fifth was not tested. Added after review: claims that deny the defining event carry about two thirds of the tenancy lead, which rests on reading such a denial as ending the process; where every fact is stated there is no advantage over the frontier prompt. On the frontier model the checklist-only setting, run once, led on tenancy and not on theft. The full product costs \sthreeFCostCasePathTenOverSameRound{} to \sthreeFCostCasePathThOverSameRound{} times the same-model prompt and holds more than needed. The tenancy references come from the engine \casepath implements, and theft reuses the conditions of an earlier study. Texts, references and rule reviews were made by AI sessions and the authors, checked by no jurist; one confirmation depends on a reordering of the hypotheses before the freeze.
\end{abstract}
